\inicializarSeccion{Investigación}
\inicializarSubseccion{Estadisticas Globales}
A nivel mundial podemos hablar del \textbf{sector industrial manufacturero} como una parte \textbf{esencial} del \textbf{desarrollo tanto económico como tecnológico y social}; actualmente se encuentra en una serie de cambios relacionados con la innovación de nuevas tecnologías.

\vspace{0.5 em}

De acuerdo a las Naciones Unidas (\citeyear{un_goal9_2026}), el valor agregado de la manufactura por persona ha aumentado 18.1 \% entre 2015 y 2025, pasando de 1,657 a 1,956 US DLLS constantes que tenía en 2020. Sin embargo, durante el mismo periodo, la participación de la manufactura en el empleo mundial disminuyó de 14.3 \% a 13.7 \%. Esto se relaciona con que las industrias están \textbf{incorporando más tecnología y procesos automatizados.} \textit{La expansión industrial se ha vuelto más intensiva en capital y tecnología, mientras que el empleo continúa moviéndose hacia otros sectores, especialmente servicios.} (\cite{un_goal9_2026})

\vspace{0.5 em}

La tecnología está cambiando el tipo de trabajo que realizan las personas en la industria, la inteligencia artificial por ejemplo y la automatización pueden sustituir algunas o varias tareas, pero también son un \textbf{complemento en el trabajo humano} para aumentar la productividad si es bien utilizada. Por eso las habilidades relacionadas con la tecnología están tomando cada vez más importancia, el manejo de equipos y la supervisión de procesos. (\cite{oit_ia_manufactura_2026})

\vspace{0.5 em}

\inicializarSubseccion{Estadisticas Nacionales}
Situándonos en México, la industria manufacturera también representa una parte muy importante de la actividad económica y el empleo. En octubre de 2025, los establecimientos manufactureros registrados tenían 2,327,216 obreros y técnicos contratados directamente por los establecimientos, lo que representa una parte importante del personal que participa directamente en las actividades productivas (\cite{inegi_immex_2025}). 

\vspace{0.5 em}

La importancia del personal que trabaja directamente en la producción se puede observar por ejemplo, la Organización Internacional del Trabajo (OIT) reportó que en México más de 1.6 millones de personas trabajaban en la industria de alimentos y bebidas en 2022, si convertimos esta cantidad en porcentaje representa un 18.4\% de los empleos manufactureros del país entero en ese sector y de acuerdo con datos del Banco Mundial (\citeyear{banco_mundial_industria_mexico_2025}), el sector industrial, incluyendo la construcción, representó aproximadamente 31.8\% del PIB de México en 2024. 

\vspace{0.5 em}

Dentro de este sector se encuentra la manufactura, que incluye actividades como la producción de automóviles, alimentos, productos electrónicos, maquinaria y diferentes bienes de consumo. Además, los obreros y técnicos de establecimientos manufactureros acumularon alrededor de 479 millones de horas trabajadas durante octubre de 2025 (\cite{inegi_immex_2025}). 

\vspace{0.5 em}

De estos datos interpretamos que aunque la industria utiliza cada vez más tecnología, sigue existiendo una gran cantidad de actividades que \textbf{requieren trabajadores capaces de operar maquinaria, utilizar herramientas, revisar procesos y mantener los equipos en funcionamiento.}

\vspace{0.5 em}

Esto se relaciona directamente con el ODS 9, porque el objetivo busca impulsar la innovación y modernizar los procesos industriales. La utilidad de nuestra caja está precisamente en aplicar una tecnología de sensores y un sistema de control a una actividad diaria del personal técnico y operativo. Así, una tarea que normalmente se hace de forma manual puede hacerse de manera más \textbf{organizada y automatizada}. La OIT también destaca que la transformación tecnológica en la manufactura requiere desarrollar nuevas habilidades y utilizar la tecnología para mejorar la productividad.

\pieDePagina{Fuente: INEGI, 2025. }

\inicializarSubseccion{Estadísticas en Nuevo León}
La importancia de la manufactura en Nuevo León se puede observar en los datos más recientes de 2026. Entre enero y agosto, el sector manufacturero generó 31,585 nuevos empleos formales en el estado. Además, en agosto se registraron 5,317 nuevos empleos manufactureros (\cite{gob_nl_empleo_septiembre_2026}). El crecimiento siguió por los siguientes meses ya que entre enero y julio de 2026, Nuevo León generó 26,268 nuevos empleos manufactureros, que representaron 64.5\% de todos los empleos creados en el estado durante ese periodo. Esto significa que aproximadamente \textbf{dos de cada tres nuevos puestos de trabajo generados en Nuevo León estuvieron relacionados con la manufactura} (\cite{gob_nl_ano_manufacturero_2026}). Para nuestro proyecto, este dato es importante porque demuestra que existe una gran cantidad de trabajadores que utilizan herramientas y equipos dentro de las industrias y que pueden beneficiarse de soluciones tecnológicas que ayuden a organizar, controlar y administrar sus herramientas de trabajo.

\vspace{0.5 em}

Nuestra caja no busca sustituir al personal técnico, sino utilizar tecnología para facilitar una parte de su trabajo. En una industria donde los procesos son cada vez más tecnológicos, tener una caja que prevea las medidas de seguridad e identifique automáticamente las herramientas mediante RFID puede ayudar a que el trabajador pierda menos tiempo buscando o revisando manualmente qué herramientas están disponibles.

\vspace{0.5 em}

\textbf{Prompt para la IA:} \textit{"Actúa como un investigador académico especializado en Industria, innovación e infraestructuras, global, nacional y local (Nuevo León Monterrey). Necesito realizar una investigación universitaria sobre Industria, innovación e infraestructuras enfocado en el ODS \#9. Construye una investigación rigurosa utilizando literatura académica y fuentes primarias siempre que sea posible. Después sintetiza las fuentes en lugar de simplemente resumirlas una por una.  Identifica también los principales debates académicos y las áreas donde existe consenso vs. desacuerdo. No inventes referencias ni datos."}

\vspace{0.5 em}

Para esta investigación se usó de base lo que generó la IA como se mencionó en la entrega y un cerebro humano para afinar y juntar la información.


\end{multicols*}